\section{Proofs for Section~\ref{sec:product}}\label{app:proofs-product}

Notation is as in Section~\ref{sec:product}. Write $\epsilon:=\epsilon_{\rm ad}$, $\mathsf V:=\mathsf V_I$ and $\mathsf c:=\mathsf c_I$ where no confusion arises. Two facts are used repeatedly.

\begin{enumerate}
\item[(P1)] For $c<V$ and $\epsilon>0$, the function $\mathcal Y_c(p)=(p-c)(1-p/V)^{\epsilon}$ is positive on $(\max\{c,0\},V)$, strictly log-concave there, strictly increasing below $p^\circ(c):=(V+\epsilon c)/(1+\epsilon)$ and strictly decreasing above it, with $\mathcal Y_c(p^\circ)=V\,\mathsf C_{\rm ad}(1-c/V)^{1+\epsilon}$. Indeed $d^2\ln\mathcal Y_c/dp^2=-(p-c)^{-2}-\epsilon(V-p)^{-2}<0$, and $d\ln\mathcal Y_c/dp=(p-c)^{-1}-\epsilon(V-p)^{-1}$ vanishes only at $p^\circ$, which lies in $(c,V)$ because $p^\circ-c=(V-c)/(1+\epsilon)>0$ and $V-p^\circ=\epsilon(V-c)/(1+\epsilon)>0$. Then $\mathcal Y_c(p^\circ)=[(V-c)/(1+\epsilon)][\epsilon(V-c)/((1+\epsilon)V)]^{\epsilon}=V\epsilon^{\epsilon}(1+\epsilon)^{-1-\epsilon}(1-c/V)^{1+\epsilon}$. If $c\ge V$, $\mathcal Y_c\le 0$ on $[0,V]$.
\item[(P2)] A strictly quasi-concave function on an interval, with peak $p^\circ$, attains its maximum over any subinterval $[a,b]$ at the point of $[a,b]$ nearest to $p^\circ$.
\end{enumerate}

\paragraph{Proof of Proposition~\ref{prop:prod-price}.}
(a) Profit per in-scope outcome at price $p\in[0,\mathsf V]$ is $(p-\hat c_I)D(p)=\mathcal Y_c(p)$ with $c=\hat c_I$, $V=\mathsf V$. If $\mathsf c<1$, (P1) gives the unique maximizer $p^m=(\mathsf V+\epsilon\hat c_I)/(1+\epsilon)$, so $\kappa^m=p^m/\mathsf V=(1+\epsilon\mathsf c)/(1+\epsilon)$, and the profit $\mathsf V\mathsf C_{\rm ad}(1-\mathsf c)^{1+\epsilon}$. The margin per outcome sold is $p^m-\hat c_I=(\mathsf V-\hat c_I)/(1+\epsilon)$, and adoption is $(1-\kappa^m)^\epsilon=[\epsilon(1-\mathsf c)/(1+\epsilon)]^\epsilon$. If $\hat c_I$ is all variable cost, the gross margin is $1-\hat c_I/p^m=1-\mathsf c/\kappa^m=(\kappa^m-\mathsf c)/\kappa^m=(1-\mathsf c)/(1+\epsilon\mathsf c)$. If $\mathsf c\ge 1$, every price earns a nonpositive profit by (P1).

(b) Let a customer of the incumbent have friction $z\mathsf V$ per outcome, with $\Pr\{z\le y\}=y^\epsilon$ for $y\in[0,1]$. Its retained saving is $\mathsf V_I-p$ from the incumbent's agent and $\mathsf V_E-(1+\upsilon_E)\hat c_E-p_{\rm sw}=\mathsf V_I-\bar p$ from an entrant's, by \eqref{eq:prod-cap}. Because frictions are common to all agents, the customer adopts an agent if and only if the larger retained saving, $\mathsf V_I-\min\{p,\bar p\}$, is at least $z\mathsf V$, and it chooses the incumbent if and only if $p\le\bar p$ (ties go to the incumbent). For $p\le\bar p$ the incumbent therefore sells to the share $\Pr\{z\le 1-p/\mathsf V\}=D(p)$, every adopter, and for $p>\bar p$ it sells nothing. Its problem is to maximize $\mathcal Y_c(p)$, $c=\hat c_I$, over $p\le\bar p$. If $\bar p\le\hat c_I$, every price that sells earns a nonpositive margin. If $\bar p>\hat c_I$, by (P1) and (P2) the maximum over $[\hat c_I,\min\{\bar p,\mathsf V\}]$ is at $\min\{p^m,\bar p\}$, because $p^m<\mathsf V$. The cap binds if and only if $\bar p<p^m$, that is, $(1+\upsilon_E)\hat c_E<p^m-\rho_u=\hat c^\ddagger_E$.

(c) Where $p^*=p^m=(\mathsf V+\epsilon\hat c_I)/(1+\epsilon)$, the derivatives are $\epsilon/(1+\epsilon)$ in $\hat c_I$ and $1/(1+\epsilon)$ in $\mathsf V$; the elasticity is $(dp^*/d\hat c_I)(\hat c_I/p^*)=\epsilon\hat c_I/(\mathsf V+\epsilon\hat c_I)$. Where $p^*=\bar p=(1+\upsilon_E)\hat c_E+\rho_u$, the price does not depend on $\hat c_I$, and $dp^*/d\hat c_E=1+\upsilon_E$. At the kink $p^m=\bar p$ these are one-sided derivatives.

(d) A vendor with $\rho_u=0$ and the entrants' cost faces the cap $(1+\upsilon_E)\hat c_E$, so this is its price whenever the cap binds. The map $(\upsilon_E,\hat c_E)\mapsto\min\{p^m,(1+\upsilon_E)\hat c_E+\rho_u\}$ is continuous, so as $\upsilon_E\to 0$ and $\hat c_E\to c_E^\infty$ the price tends to $\min\{p^m,c_E^\infty+\rho_u\}$, and where the cap binds the capture share tends to $(c_E^\infty+\rho_u)/\mathsf V$. If also $\hat c_I\to c_E^\infty$, the margin per dollar of revenue, $(p^*-\hat c_I)/p^*$, tends to $\rho_u/(c_E^\infty+\rho_u)$ where the cap binds; with $\rho_u=0$ it tends to zero.

(e) By definition, $\mathsf V_I-\mathsf V_E=c^{\rm res}_E-c^{\rm res}_I=[c^Q(e_E)-c^Q(e_I)]+(\eta_E-\eta_I)$, which gives the decomposition. By Proposition~\ref{prop:tasks-review}, $c^Q(e)=\min_{\xi}\{wt^R\xi+eL[1-\xi(1+\ln(1/\xi))]\}$ over the admissible review shares, and by the envelope theorem its slope in $e$ is $L[1-\xi^*(1+\ln(1/\xi^*))]$ at the optimal share $\xi^*$, whether the economic term or the floor binds. The function $\xi\mapsto\xi(1+\ln(1/\xi))$ maps $(0,1)$ into $(0,1)$, so $0<dc^Q/de<L$. Since $e(\mathcal D)$ is strictly decreasing, the data rent $c^Q(e(\mathcal D_E))-c^Q(e(\mathcal D_I))$ is strictly increasing in $\mathcal D_I$ and, by the mean value theorem, at most $L[e(\mathcal D_E)-e(\mathcal D_I)]=La_e(\mathcal D_E^{-b_e}-\mathcal D_I^{-b_e})$ for $\mathcal D_I>\mathcal D_E$. With $\mathcal D_I=\lambda\mathcal D_E$ for a fixed $\lambda>1$, the bound is $La_e\mathcal D_E^{-b_e}(1-\lambda^{-b_e})\to 0$ as $\mathcal D_E\to\infty$, and the rent is nonnegative, so it tends to zero. Where the floor binds, $\xi^*=\xi^{\min}$ is fixed and $c^Q$ is affine and increasing in $e$; the rent is then a constant minus a positive multiple of $e(\mathcal D_I)$, and $e''(\mathcal D)=a_eb_e(1+b_e)\mathcal D^{-b_e-2}>0$, so the rent is concave in $\mathcal D_I$.

(f) A customer with premium $p_{\rm sw}$ keeps the incumbent's agent if and only if $\mathsf V_I-p\ge\mathsf V_E-(1+\upsilon_E)\hat c_E-p_{\rm sw}=\mathsf V_I-\bar p_0-p_{\rm sw}$, that is, $p_{\rm sw}\ge p-\bar p_0$ (ties go to the incumbent), and then adopts if and only if $\mathsf V_I-p\ge z\mathsf V$. By independence the incumbent's sales are $D(p)\mathsf S_{\rm sw}(p-\bar p_0)$ and its profit is $\Pi(p):=\mathcal Y_c(p)\mathsf S_{\rm sw}(p-\bar p_0)$, $c=\hat c_I$. $\mathsf S_{\rm sw}$ is nonincreasing and upper semicontinuous and $\mathcal Y_c$ is continuous, so $\Pi$ attains its maximum on $[\hat c_I,\mathsf V_I]$. By (P1), $\mathcal Y_c$ is strictly increasing below $p^m$ and strictly decreasing above it. For $p\le\bar p_0$, $\Pi=\mathcal Y_c$. If $p^m<\bar p_0$, then $\Pi(p)\le\mathcal Y_c(p)<\mathcal Y_c(p^m)=\Pi(p^m)$ for $p\ne p^m$, so $p^*=p^m$. If $p^m\ge\bar p_0$, then $\Pi(p)=\mathcal Y_c(p)<\mathcal Y_c(\bar p_0)=\Pi(\bar p_0)$ for $p<\bar p_0$, and $\Pi(p)\le\mathcal Y_c(p)\mathsf S_{\rm sw}(p^m-\bar p_0)<\Pi(p^m)$ for $p>p^m$ whenever $\Pi(p^m)>0$ (if $\Pi(p^m)=0$, then $\Pi=0$ above $p^m$ and a maximizer lies in $[\bar p_0,p^m]$). Hence a maximizer lies in $[\min\{p^m,\bar p_0\},p^m]$. With $\mathsf S_{\rm sw}(x)=\mathbf 1\{x\le p_{\rm sw}\}$, $\Pi(p)=\mathcal Y_c(p)$ for $p\le\bar p_0+p_{\rm sw}=\bar p$ and zero above, which is the problem of (b).

With a log-concave density, $\mathsf S_{\rm sw}$ is log-concave, continuous, equal to 1 at zero and differentiable, with nondecreasing hazard $h_{\rm sw}$ (assumed differentiable below). On $(\bar p_0,\mathsf V_I)$, where $\bar p_0>\hat c_I$, $\ln\Pi=\ln(p-\hat c_I)+\epsilon_{\rm ad}\ln(1-p/\mathsf V_I)+\ln\mathsf S_{\rm sw}(p-\bar p_0)$ is strictly concave, with derivative $\ell(p)-h_{\rm sw}(p-\bar p_0)$, and $\Pi$ is continuous at $\bar p_0$. Since $\Pi=\mathcal Y_c$ is increasing on $[\hat c_I,\bar p_0]$ (because $\bar p_0<p^m$), the maximizer is $\bar p_0$ if the right derivative $\ell(\bar p_0)-h_{\rm sw}(0)\le 0$, and otherwise the unique root of \eqref{eq:prod-het} on $(\bar p_0,\mathsf V_I)$, which is below $p^m$ because $\ell\le 0<h_{\rm sw}$ on $[p^m,\mathsf V_I)$ wherever the density is positive. Write \eqref{eq:prod-het} as $G(p):=\ell(p)-h_{\rm sw}(p-\bar p_0)=0$. Then $-\partial G/\partial p=\mathcal G:=(p-\hat c_I)^{-2}+\epsilon_{\rm ad}(\mathsf V_I-p)^{-2}+h'_{\rm sw}>0$, $\partial G/\partial\hat c_I=(p-\hat c_I)^{-2}$ and $\partial G/\partial\bar p_0=h'_{\rm sw}\ge 0$, so by the implicit function theorem
\[
\frac{dp^*}{d\hat c_I}=\frac{(p-\hat c_I)^{-2}}{\mathcal G}\in(0,1),\qquad
\frac{dp^*}{d\hat c_E}=(1+\upsilon_E)\frac{h'_{\rm sw}}{\mathcal G}\in[0,1+\upsilon_E).
\]
Exponential premia have $h_{\rm sw}\equiv 1/\bar p_{\rm sw}$ and $h'_{\rm sw}=0$, which gives the threshold and $dp^*/d\hat c_E=0$; $\ell(\bar p_0)>0$ because $\bar p_0<p^m$.

Gamma premia with mean $\bar p_{\rm sw}$ and coefficient of variation $\mathrm{cv}_{\rm sw}$ have shape $\mathrm{cv}_{\rm sw}^{-2}\ge 1$, hence a log-concave density, and variance $(\mathrm{cv}_{\rm sw}\bar p_{\rm sw})^2$. Let $\bar p:=\bar p_0+\bar p_{\rm sw}>\hat c_I$, let $p_n$ maximize $\Pi_n$ along a sequence $\mathrm{cv}_n\to 0$, and let $Y^*:=\max_{\hat c_I\le p\le\bar p}\mathcal Y_c(p)=\mathcal Y_c(\min\{p^m,\bar p\})>0$. By Chebyshev's inequality $\mathsf S_n(x)\to 1$ uniformly on $x\le\bar p_{\rm sw}-\delta$ and $\mathsf S_n(x)\to 0$ uniformly on $x\ge\bar p_{\rm sw}+\delta$, for every $\delta>0$. Hence $\Pi_n(p_n)\ge\Pi_n(\min\{p^m,\bar p-\delta\})\to\mathcal Y_c(\min\{p^m,\bar p-\delta\})$, which tends to $Y^*$ as $\delta\to0$, while $\sup_{p\ge\bar p+\delta}\Pi_n(p)\le\max\mathcal Y_c\cdot\mathsf S_n(\bar p_{\rm sw}+\delta)\to 0$. So every limit point $\hat p$ of $(p_n)$ satisfies $\hat p\le\bar p$ and, since $\Pi_n\le\mathcal Y_c$ and $\mathcal Y_c$ is continuous, $\mathcal Y_c(\hat p)\ge Y^*$; as $\mathcal Y_c$ has a unique maximizer on $[\hat c_I,\bar p]$, $\hat p=\min\{p^m,\bar p\}$, the price of (b). \qed

\paragraph{Proof of Proposition~\ref{prop:prod-seats}.}
(a) At constant volume, seats are $N_{\rm seat}(0)[1-\epsilon_{\rm seat}\mathcal X_h(t)]=N_{\rm seat}(0)[1-\mathcal X(t)]$ by Assumption~\ref{ass:prod-seats}; with every activity's volume growing at $g_v$, they are $N_{\rm seat}(0)e^{g_vt}[1-\mathcal X(t)]$, and revenue is $p_{\rm seat}$ times this. Its log-derivative is $g_v+d\ln(1-\mathcal X)/dt=g_v-\nu$. At constant volume, the revenue lost against the uncompressed base is $p_{\rm seat}N_{\rm seat}(0)\mathcal X(t)=p_{\rm seat}N_{\rm seat}(0)\epsilon_{\rm seat}\mathcal X_h(t)$. If $\mathcal X$ depends on time only through the horizon, with $\log_2H(t)=\log_2H_0+(t-t_0)/T_d$ (Assumption~\ref{ass:ext-capability}(i)), then $d\ln N_{\rm seat}/d\log_2H=(d\ln N_{\rm seat}/dt)\,T_d=-T_d\nu$.

(b) Let $\omega^{\rm seat}_j=v_jt^{\rm seat}_{x^0_j}/\sum_iv_it^{\rm seat}_{x^0_i}$. With the share $\mathsf B(t-T_j)$ of activity $j$'s volume moved to $x^1_j$ by date $t$ (with $\mathsf B(s)=0$ for $s<0$), seat-hours relative to their initial level are $\sum_j\omega^{\rm seat}_j[1-\mathsf B(t-T_j)+\mathsf B(t-T_j)t^{\rm seat}_{x^1_j}/t^{\rm seat}_{x^0_j}]$, which is $1-\mathcal X_h(t)$ for the stated $\mathcal X_h$.

(c) Differentiating, $\mathsf B$ solves $\dot{\mathsf B}=(\mathfrak p+\mathfrak q\mathsf B)(1-\mathsf B)$ with $\mathsf B(0)=0$ (checked symbolically by the script). Hence $\ddot{\mathsf B}=\dot{\mathsf B}[\mathfrak q(1-\mathsf B)-(\mathfrak p+\mathfrak q\mathsf B)]=\dot{\mathsf B}(\mathfrak q-\mathfrak p-2\mathfrak q\mathsf B)$. Since $\dot{\mathsf B}>0$ and $\mathsf B$ rises from 0 to 1, $\dot{\mathsf B}$ increases while $\mathsf B<(\mathfrak q-\mathfrak p)/(2\mathfrak q)$ and decreases afterwards, so if $\mathfrak q>\mathfrak p$ it peaks when $\mathsf B=(\mathfrak q-\mathfrak p)/(2\mathfrak q)$, where $\dot{\mathsf B}=\tfrac{\mathfrak p+\mathfrak q}{2}\cdot\tfrac{\mathfrak p+\mathfrak q}{2\mathfrak q}=(\mathfrak p+\mathfrak q)^2/(4\mathfrak q)$. Inverting $\mathsf B$, $e^{-(\mathfrak p+\mathfrak q)t}=(1-\mathsf B)/(1+\mathfrak q\mathsf B/\mathfrak p)$; at $\mathsf B=(\mathfrak q-\mathfrak p)/(2\mathfrak q)$ this ratio is $\mathfrak p/\mathfrak q$, so $t^*_B=\ln(\mathfrak q/\mathfrak p)/(\mathfrak p+\mathfrak q)$. Multiplying by $\bar x$ gives the rate for $\mathcal X$.

For the seat-loss rate, $\nu=\bar x\dot{\mathsf B}/(1-\bar x\mathsf B)$ and $\nu(0)=\bar x\mathfrak p$. Its log-derivative is $\ddot{\mathsf B}/\dot{\mathsf B}+\bar x\dot{\mathsf B}/(1-\bar x\mathsf B)=\psi(\mathsf B(t))$, with
\[
\psi(F):=\mathfrak q-\mathfrak p-2\mathsf qF+\frac{\bar x(\mathfrak p+\mathsf qF)(1-F)}{1-\bar xF},\qquad F\in[0,1].
\]
Write $N(F)=(\mathfrak p+\mathsf qF)(1-F)$ and $\Delta_F=1-\bar xF>0$. Then $N'=\mathfrak q-\mathfrak p-2\mathsf qF$, and a direct expansion gives $N'\Delta_F+\bar xN=\mathfrak q(1-2F+\bar xF^2)-\mathfrak p(1-\bar x)$, so that
\[
-\psi'(F)\Delta_F^2=2\mathfrak q\Delta_F^2-\bar x\bigl[\mathfrak q(1-2F+\bar xF^2)-\mathfrak p(1-\bar x)\bigr]=\mathfrak q\bigl[\Delta_F^2+1-\bar x\bigr]+\bar x\mathfrak p(1-\bar x)>0 .
\]
So $\psi$ is strictly decreasing, and since $\mathsf B$ is strictly increasing, $d\ln\nu/dt$ is strictly decreasing in $t$, so $\nu$ is single-peaked or monotone. Now $\psi(0)=\mathfrak q-\mathfrak p(1-\bar x)$ and, for $\bar x<1$, $\psi(1)=-(\mathfrak p+\mathfrak q)<0$. If $\mathfrak q>\mathfrak p(1-\bar x)$, $\nu$ rises until $\psi(\mathsf B)=0$ and falls afterwards. Since $\psi(F)\Delta_F=\mathfrak q(1-2F+\bar xF^2)-\mathfrak p(1-\bar x)$, the root solves $\bar x\mathsf qF^2-2\mathsf qF+\mathfrak q-\mathfrak p(1-\bar x)=0$, whose roots are $[1\pm\sqrt{(1-\bar x)(1+\bar x\mathfrak p/\mathfrak q)}]/\bar x$. The larger exceeds $1/\bar x>1$. The smaller is \eqref{eq:prod-Bhat}; it is positive because $(1-\bar x)(1+\bar x\mathfrak p/\mathfrak q)<1$ is equivalent to $\mathfrak q>\mathfrak p(1-\bar x)$, and below one because $(1-\bar x)(1+\bar x\mathfrak p/\mathfrak q)>(1-\bar x)^2$. If $\mathfrak q\le\mathfrak p(1-\bar x)$, $\psi<0$ for $t>0$ and $\nu$ falls from the start. As $t\to\infty$, $\dot{\mathsf B}\to 0$ while $1-\bar x\mathsf B\to 1-\bar x>0$, so $\nu\to 0$. At $t^*_B$, $\ddot{\mathsf B}=0$ and $d\ln\nu/dt=\bar x\dot{\mathsf B}/(1-\bar x\mathsf B)>0$, so the peak of $\nu$ comes after $t^*_B$.

(d) $\dot{\mathcal X}=\bar x\,\mathsf n((t-t_m)/s_T)/s_T$ with $s_T=\varsigma_{\rm dbl}T_d$, maximal at $t=t_m$, where $\mathsf n(0)=1/\sqrt{2\pi}$. The form of $\mathcal X$ is the probit share of Proposition~\ref{prop:ext-feasible}(d). \qed

\paragraph{Proof of Proposition~\ref{prop:prod-migrate}.}
(a) Write $R_S(t)=R_0e^{g_vt}[1-\mathcal X(t)]$ and $r_v:=r-g_v>0$. Then
\[
V_k(T)=\varpi_k\gamma[I_1(T)-I_2(T)]-e^{-rT}\mathcal T_k(T),
\]
with
\[
\begin{aligned}
I_1(T)&=(1-\delta_k)R_0\int_T^\infty e^{-r_vt}\bigl[(1-e_k)(1-\mathcal X(T))+e_k(1-\mathcal X(t))\bigr]dt,\\
I_2(T)&=R_0\int_T^\infty e^{-r_vt}(1-\mathcal X(t))\,dt .
\end{aligned}
\]
The integrands are bounded and $r_v>0$, so both integrals converge and can be differentiated under the integral sign, with $\mathcal X$ continuously differentiable:
\[
\begin{aligned}
I_1'(T)&=-(1-\delta_k)R_0e^{-r_vT}(1-\mathcal X(T))-(1-\delta_k)(1-e_k)R_0\dot{\mathcal X}(T)\frac{e^{-r_vT}}{r_v},\\
I_2'(T)&=-R_0e^{-r_vT}(1-\mathcal X(T)).
\end{aligned}
\]
Hence
\[
I_1'-I_2'=R_0e^{-r_vT}(1-\mathcal X(T))\Bigl[\delta_k-\frac{(1-\delta_k)(1-e_k)\nu(T)}{r_v}\Bigr]=e^{-rT}R_S(T)\Bigl[\delta_k-\frac{(1-\delta_k)(1-e_k)\nu(T)}{r_v}\Bigr].
\] Also $d[-e^{-rT}\mathcal T_k(T)]/dT=e^{-rT}[r\mathcal T_k-\mathcal T_k']$, and since $\dot R_S=R_S(g_v-\nu)$, $r\mathcal T_k-\mathcal T_k'=\omega_k\gamma R_S(T)[r-g_v+\nu(T)]$. Adding the two parts gives (a).

(b) The sign of $V_k'(T)$ is that of $A(\nu):=\varpi_k\delta_k+\omega_kr_v-\nu c_1$, with $c_1:=\varpi_k(1-\delta_k)(1-e_k)/r_v-\omega_k$. If $c_1>0$, $A(\nu)<0$ if and only if $\nu>[\varpi_k\delta_k+\omega_kr_v]/c_1=\nu^*_k$, which is \eqref{eq:prod-nustar}. If $c_1\le 0$, then $A\ge\varpi_k\delta_k+\omega_kr_v\ge 0$ for all $\nu\ge 0$, so $V_k$ is nondecreasing; since both parts of $V_k(T)$ tend to zero as $T\to\infty$, $V_k\le 0$ everywhere. For the comparative statics write $\nu^*_k=r_vN/M$ with $N=\varpi_k\delta_k+\omega_kr_v$ and $M=\varpi_k(1-\delta_k)(1-e_k)-\omega_kr_v>0$. Raising $\delta_k$ raises $N$ and lowers $M$; raising $\omega_k$ raises $N$ and lowers $M$; raising $e_k$ lowers $M$ and leaves $N>0$ unchanged. With $a:=(1-\delta_k)(1-e_k)$, $\partial(N/M)/\partial\varpi_k$ has the sign of $\delta_kM-aN=-\omega_kr_v(\delta_k+a)$, which is negative when $\omega_k>0$.

(c) If $\nu$ is single-peaked or decreasing, the set $\{t:\nu(t)>\nu^*_k\}$ is an interval $(T_1,T_2)$, possibly empty, with $T_1=0$ if $\nu(0)\ge\nu^*_k$. By (b), $V_k$ is nondecreasing on $[0,T_1]$, decreasing on $[T_1,T_2]$ and nondecreasing on $[T_2,\infty)$, where it rises to its limit $0$. So $\sup_TV_k(T)=\max\{V_k(T_1),0\}$. If $V_k(T_1)>0$, $T_1$ is optimal; otherwise no date gives a positive value. If the interval is empty, $V_k$ is nondecreasing and migration never pays.

(d) For each $T$: $\partial V_k/\partial\delta_k=-\varpi_k\gamma(1-\delta_k)^{-1}\int_T^\infty e^{-rt}R_k(t)dt<0$; $\partial V_k/\partial\omega_k=-e^{-rT}\gamma R_S(T)<0$; $R_k$ depends on $e_k$ through $1-e_k+e_k(1-\mathcal X(t))/(1-\mathcal X(T))$, whose derivative in $e_k$ is $(1-\mathcal X(t))/(1-\mathcal X(T))-1\le 0$ because $\mathcal X$ is nondecreasing; and $\partial V_k/\partial\varpi_k=[V_k(T)+e^{-rT}\mathcal T_k(T)]/\varpi_k$, which is positive where $V_k(T)>0$. The maximum over $T$ of functions that are monotone in a parameter is monotone in the same direction. \qed

\paragraph{Proof of Proposition~\ref{prop:prod-portfolio}.}
(a) In the small-share limit of Proposition~\ref{prop:dynamic}(d), revenue grows at $g^\circ$, the firm's weight is $v_0$, and Proposition~\ref{prop:ext-customer}(b) (with zero-spend growth $G_i$ in place of $G$ and the signs reversed for gains) gives, for a permanent change, $dv_0=[dm_0+(v_0-k)(dG_i+\Delta^\circ d\lambda/\lambda)]/(r-g^\circ)$, exact to first order. A feature at full adoption raises gross profit per dollar of base revenue by $\upsilon_{\mathsf f}-\pi n_{\mathsf f}$, a change in $m_0$, and raises $G_i$ by $\delta^G_{\mathsf f}$ and $\lambda$ by the proportion $\delta^\lambda_{\mathsf f}$. A time-varying perturbation is a superposition of permanent perturbations that start at different dates. The small-share problem is stationary, so a perturbation that starts at $\tau$ is worth $e^{-(r-g^\circ)\tau}$ times the same perturbation starting at 0, and first-order effects add. Hence value changes by $S\int_0^\infty e^{-(r-g^\circ)t}R_{\mathsf f}(t)dt$, to which the development cost $F_{\mathsf f}$ is added.

(b) Without capacity limits, the objective is $\sum_{\mathsf f}x_{\mathsf f}\mathrm{NPV}_{\mathsf f}-F^{\rm plat}\mathbf 1_{\rm plat}$, where $\mathbf 1_{\rm plat}=1$ if some feature that needs the platform is built. Given the platform, each feature that needs it should be built if and only if its NPV is positive, which yields $\sum(\mathrm{NPV}_{\mathsf f})^+-F^{\rm plat}$; without it, those features yield nothing. Features that do not need the platform are chosen separately. With capacity, taking the platform decision as given, the relaxation $\max\sum x_{\mathsf f}\mathrm{NPV}_{\mathsf f}$ subject to $\sum x_{\mathsf f}n^E_{\mathsf f}\le N^E$, $0\le x_{\mathsf f}\le 1$, is a fractional knapsack. Let $\mu$ be the ratio $\mathrm{NPV}_{\mathsf f}/n^E_{\mathsf f}$ of the feature at which capacity runs out (or $\mu=0$ if it does not). The greedy solution and $\mu$ satisfy the optimality conditions of the linear program: features with ratio above $\mu$ are at $x=1$ and have reduced cost $\mathrm{NPV}_{\mathsf f}-\mu n^E_{\mathsf f}\ge 0$; those with ratio below $\mu$ are at $x=0$; capacity binds if $\mu>0$. So the greedy solution is optimal and $\mu$ is the dual price of capacity, the derivative of the relaxed value in $N^E$ where it is differentiable. Dropping the fractional feature gives a feasible integer portfolio, so the integer optimum lies between the relaxed value less one feature's NPV and the relaxed value.

(c) At full adoption $R_{\mathsf f}=\upsilon_{\mathsf f}-\pi n_{\mathsf f}+(X-k)(\delta^G_{\mathsf f}+\Delta^*\delta^\lambda_{\mathsf f})$, which is positive if and only if the stated inequality holds. Replacing $v_0$ by $X^{\rm scr}$ multiplies the term $(X-k)(\cdot)$ by $(X^{\rm scr}-k)/(v_0-k)$. \qed

\paragraph{Proof of Proposition~\ref{prop:prod-model}.}
(a) Under a binding cap, option $\mathsf m$ sells $Q$ outcomes at $(1+\upsilon_E)\hat c_E+\rho_{u,\mathsf m}$ and bears $Q[c^{\rm var}_{\mathsf m}+(r+\delta_K)K_{\mathsf m}]+\tilde F_{\mathsf m}$, so its gross profit is $Q[(1+\upsilon_E)\hat c_E-\hat c^{\rm net}_{\mathsf m}]-\tilde F_{\mathsf m}$. The entrants' term is common to all options, so the optimal option minimizes $C_{\mathsf m}(Q):=Q\hat c^{\rm net}_{\mathsf m}+\tilde F_{\mathsf m}$. If $\mathsf m$ is optimal at $Q$ and $\mathsf m'$ at $Q'>Q$, then $C_{\mathsf m}(Q)\le C_{\mathsf m'}(Q)$ and $C_{\mathsf m'}(Q')\le C_{\mathsf m}(Q')$; adding, $(Q'-Q)(\hat c^{\rm net}_{\mathsf m'}-\hat c^{\rm net}_{\mathsf m})\le 0$, so the net unit cost of the optimal option is nonincreasing in $Q$, and its fixed cost is then nondecreasing. Between two neighbors on the lower envelope the switch occurs where $C_{\mathsf m}=C_{\mathsf m'}$, which is $Q_{\mathsf m\mathsf m'}$.

(b) $\hat c^{\rm net}_{\mathrm{api}}(t)=\pi_{\mathrm{api}}(0)e^{-\delta_\pi t}n_{\mathrm{api}}+c^h_{\mathrm{api}}-\rho_{u,\mathrm{api}}$ is decreasing in $t$, and the other net unit costs are constant. For an option with $\tilde F_{\mathsf m}>\tilde F_{\mathrm{api}}$, $Q_{\mathrm{api},\mathsf m}(t)=(\tilde F_{\mathsf m}-\tilde F_{\mathrm{api}})/[\hat c^{\rm net}_{\mathrm{api}}(t)-\hat c^{\rm net}_{\mathsf m}]$ while the denominator is positive and $+\infty$ afterwards; it is nondecreasing because the denominator decreases. In the limit, $\hat c^{\rm net}_{\mathrm{api}}\to c^h_{\mathrm{api}}-\rho_{u,\mathrm{api}}$. If $\hat c^{\rm net}_{\mathsf m}\ge c^h_{\mathrm{api}}-\rho_{u,\mathrm{api}}$, then eventually the API has a weakly lower net unit cost and a lower fixed cost, so it is weakly better at every volume; hence option $\mathsf m$ can be optimal at some volume in the long run only if $\hat c^{\rm net}_{\mathsf m}<c^h_{\mathrm{api}}-\rho_{u,\mathrm{api}}$, which rearranges to the stated condition.

(c) Capital per dollar of revenue rises by $K_{\mathsf m}Q/S$; Proposition~\ref{prop:tasks-capital}(a)--(c) applies with $\Delta K_j=K_{\mathsf m}$ and $v_j=Q$.

(d) The partnership has unit cost $\kappa_{\rm par}p+c^h$ and no fixed cost, so it is the option with the lowest fixed cost in (a), optimal below the first threshold on the envelope; a partner input that raises $\rho_u$ lowers its net unit cost, and one that raises $\lambda$ adds $(X-k)\Delta^*\Delta\lambda/\lambda$ to the score by Proposition~\ref{prop:envelope}. \qed

\paragraph{Proof of Proposition~\ref{prop:prod-margin}.}
(a) $\gamma_{\rm AI}=1-f-u/p$ and $u/p=(u_0/p_0)(u/u_0)^{1-\varphi}$, with $u_0/p_0=1-f-\gamma_0$. Differentiating, $\dot\gamma_{\rm AI}=-(1-f-\gamma_0)(1-\varphi)(u/u_0)^{-\varphi}\dot u/u_0$, and $\dot u=-\delta_uu_ce^{-\delta_ut}<0$; with $\gamma_0<1-f$, the derivative is positive if and only if $\varphi<1$. As $t\to\infty$, $u\to u_h$, which gives $\gamma_\infty$.

(b) For $\varphi<1$, $\gamma_{\rm AI}$ is continuous and strictly increasing toward $\gamma_\infty$, so it reaches $\gamma_{\rm sub}$ in finite time if and only if $\gamma_\infty>\gamma_{\rm sub}$. Setting $\gamma_{\rm AI}=\gamma_{\rm sub}$ gives $(u/u_0)^{1-\varphi}=z$, so $u_h+u_ce^{-\delta_uT_\times}=u_0z^{1/(1-\varphi)}$, which is \eqref{eq:prod-Tx}. The right side exceeds $u_h$ exactly when $\gamma_\infty>\gamma_{\rm sub}$, and is below $u_0$ because $z<1$, so $T_\times>0$. With $u_h=0$, $e^{-(1-\varphi)\delta_uT_\times}=z$. For $\varphi\ge 1$ the margin never rises.

(c) With $p=u+\rho_u$, $\gamma_{\rm AI}=1-u/(u+\rho_u)=\rho_u/(u+\rho_u)$, whose derivative in $u$ is $-\rho_u/(u+\rho_u)^2<0$; the elasticity of $p$ in $u$ is $u/(u+\rho_u)$. As $u\to u_h$, the margin tends to $\rho_u/(u_h+\rho_u)$, which exceeds $\gamma_{\rm sub}$ if and only if $\rho_u/u_h>\gamma_{\rm sub}/(1-\gamma_{\rm sub})$. With $\rho_u=0$ the margin is zero.

(d) $d[w\gamma_{\rm AI}+(1-w)\gamma_{\rm sub}]/dt=\dot w(\gamma_{\rm AI}-\gamma_{\rm sub})+w\dot\gamma_{\rm AI}$, with $\gamma_{\rm sub}$ constant. \qed

\paragraph{Proof of Proposition~\ref{prop:prod-cannibal}.}
(a) Without a launch, all in-scope work is done by people and earns seat gross profit $\gamma_{\rm seat}\bar p_{\rm seat}$ per unit. With a launch at $p$, the share $D(p)$ is automated: it pays $p-\hat c_I$ and no longer earns seat profit. The change is $(p-\hat c_I-\gamma_{\rm seat}\bar p_{\rm seat})D(p)$, which is $\mathcal Y_c(p)$ of (P1) with $c=\hat c_I+\gamma_{\rm seat}\bar p_{\rm seat}=(\mathsf c+\mathsf s)\mathsf V$. By (P1), launching pays if and only if $\mathsf c+\mathsf s<1$, at the capture share $[1+\epsilon(\mathsf c+\mathsf s)]/(1+\epsilon)$, with gain $\mathsf V\mathsf C_{\rm ad}(1-\mathsf c-\mathsf s)^{1+\epsilon}$. Both comparisons follow from $\mathsf s>0$.

(b) If the incumbent does not launch, its customers adopt entrants' agents when their retained saving, $\mathsf V_I-\bar p$ by the proof of Proposition~\ref{prop:prod-price}(b), exceeds their friction, so the share $D(\bar p)$ is automated and seat profit falls to $\gamma_{\rm seat}\bar p_{\rm seat}[1-D(\bar p)]$. If it launches at $p\le\bar p$, it earns $\gamma_{\rm seat}\bar p_{\rm seat}[1-D(p)]+(p-\hat c_I)D(p)$. At $p=\bar p$ the difference is $(\bar p-\hat c_I)D(\bar p)$. (Over $p\le\bar p$ the objective is $\gamma_{\rm seat}\bar p_{\rm seat}+(p-\hat c_I-\gamma_{\rm seat}\bar p_{\rm seat})D(p)$, maximized at $\min\{\kappa^m_{\rm can}\mathsf V,\bar p\}$ by (P1) and (P2): pricing below the cap would add cannibalization on marginal adopters only.)

(c) Measure flows per unit of date-0 work, which grows at $g_v$, and let $r_v=r-g_v$ and $\tau\sim\mathrm{Exp}(\nu_E)$ be the entry date. Relative to waiting and launching at entry, launching now at $p_L$ changes flows by $b:=(p_L-\hat c_I-\gamma_{\rm seat}\bar p_{\rm seat})D(p_L)$ before $\tau$ and by $a:=p_{\rm sw}D(p_L)$ after it. The expected present value of the difference is
\[
\mathbb E\Bigl[\int_0^\tau e^{-r_vt}b\,dt+\int_\tau^\infty e^{-r_vt}a\,dt\Bigr]=\frac b{r_v}\Bigl(1-\frac{\nu_E}{r_v+\nu_E}\Bigr)+\frac a{r_v}\cdot\frac{\nu_E}{r_v+\nu_E}=\frac{b+\nu_Ea/r_v}{r_v+\nu_E}=\frac{(p_L-c_{\rm eff})D(p_L)}{r_v+\nu_E},
\]
using $\mathbb Ee^{-r_v\tau}=\nu_E/(r_v+\nu_E)$. A plan to launch at a date $T>0$ if entry has not occurred, and at entry otherwise, differs from waiting only on the event $\{\tau>T\}$, which has probability $e^{-\nu_ET}$; on that event, by the memorylessness of $\tau$ and the constant flows per unit of work, the continuation is the date-0 problem scaled by $e^{-rT}e^{g_vT}$. Its value is therefore $e^{-(r_v+\nu_E)T}$ times that of launching now. So the incumbent launches now if the latter is positive for some $p_L\ge 0$, and waits for entry otherwise. Since $D(p_L)>0$ requires $p_L<\mathsf V$, a positive value exists if and only if $c_{\rm eff}<\mathsf V$. Writing out $c_{\rm eff}$, this holds whenever $\hat c_I+\gamma_{\rm seat}\bar p_{\rm seat}<\mathsf V$, and otherwise if and only if $\nu_E>\nu_E^*$. The value $(p_L-c_{\rm eff})D(p_L)$ is $\mathcal Y_c(p_L)$ of (P1) with $c=c_{\rm eff}$; its unconstrained maximizer is $(\mathsf V+\epsilon c_{\rm eff})/(1+\epsilon)$, and with $p_L\ge 0$, (P2) gives $p_L=\max\{0,(\mathsf V+\epsilon c_{\rm eff})/(1+\epsilon)\}$. If $\nu_E>0$, $c_{\rm eff}<\hat c_I+\gamma_{\rm seat}\bar p_{\rm seat}$, so $p_L<\kappa^m_{\rm can}\mathsf V$. \qed

\paragraph{Proof of Proposition~\ref{prop:prod-aggregate}.}
(a) Gross profit is additive across lines, $\gamma S=\sum_l\gamma_lS_l$. Revenue is additive, so $\dot S/S=\sum_l(S_l/S)\dot S_l/S_l$; applying this to the parts of growth that need no acquisition spend gives $G_i=\sum_lw_lG_l$. The seat-line rate is Proposition~\ref{prop:prod-seats}(a). A tariff that bills per outcome handled on the platform is indexed to work volume, which grows at $g_v$. An outcome-priced agent's revenue is $p\,D(p)\,\mathsf B(t)\,\bar x\,ve^{g_vt}$ for in-scope volume $ve^{g_vt}$, whose log-derivative is the stated rate. Feature effects are defined per unit of adoption; to first order, the proportional changes in $\lambda$ multiply, so they add. Under Remark~\ref{rem:assumptions}(c), $k=k_{\rm fix}-k_{\rm def}$, so removing deferred revenue raises $k$ by the amount removed.

(b) With in-scope work $v$, seat revenue on the uncompressed base is $M_{\rm seat}=\bar p_{\rm seat}v$. After customers automate the share $\mathcal X$, seat revenue is $M_{\rm seat}(1-\mathcal X)$ and agent revenue is $\kappa_{\rm cap}\mathsf V\mathcal Xv=M_{\rm seat}\mathcal X\kappa_{\rm cap}\mathsf W$, since $\mathsf W=\mathsf V/\bar p_{\rm seat}$. Their sum is $M_{\rm seat}[1+\mathcal X(\kappa_{\rm cap}\mathsf W-1)]$, which exceeds $M_{\rm seat}$ for $\mathcal X>0$ if and only if $\kappa_{\rm cap}\mathsf W>1$. Its log-derivative, holding $M_{\rm seat}$'s own growth fixed, is the stated term.

(c) A seat price fixed for the contract term does not respond to unit cost, so $\varphi=0$; under a binding cap, $dp^*/d\hat c_I=0$ by Proposition~\ref{prop:prod-price}(c); the monopoly elasticity is part (c) of that proposition; and a cost-plus price $p=(1+\upsilon)u$ has $d\ln p/d\ln u=1$. Lemma~\ref{lem:prices}(a) gives $d\gamma=-(1-\varphi)du/p$, which is the stated effect per dollar of revenue (Proposition~\ref{prop:tasks-aggregate}(a)).

(d) With zero-spend growth $G_i$ (Remark~\ref{rem:assumptions}(b)) and the depreciation adjustment of Remark~\ref{rem:assumptions}(c), \eqref{eq:Rstar} reads $R^*_X=(\gamma-a-\delta_Kk)+(X-k)G_i+C_\alpha[\lambda h(X-k)]^{1/(1-\alpha)}$. It is linear in $\gamma$ and $G_i$, and the other derivatives are those of Proposition~\ref{prop:envelope} with $\partial R^*_X/\partial k=-(g^*+\delta_K)$, which gives \eqref{eq:prod-dR}. The value statement is Proposition~\ref{prop:tasks-aggregate}(c). For a transient term, $\int_0^\infty\Delta G_i(t)\,dt=\lim_{t\to\infty}\ln[S(t)/S^0(t)]$, the log of the revenue gain against the baseline path $S^0$; for the market term this is $\ln[1+\mathcal X_\infty(\kappa_{\rm cap}\mathsf W-1)]$ by (b).

(e) The objective is $U(\varphi)+(1-\varphi)|\Delta u(z)|-C(z)$. For $z'>z$, the difference in the objective between $z'$ and $z$ is $(1-\varphi)[|\Delta u(z')|-|\Delta u(z)|]-[C(z')-C(z)]$, which is nondecreasing in $1-\varphi$ because the saving is nondecreasing in $z$. The objective thus has increasing differences in $(1-\varphi,z)$, and by the monotone selection theorem \citep{milgrom1994monotone} the largest maximizer in $z$ is nondecreasing in $1-\varphi$. \qed

\paragraph{Proof of Corollary~\ref{cor:prod-dilution}.}
(i) By the proof of Proposition~\ref{prop:prod-aggregate}(d), $R^*_X$ depends on $m_0$ and $G_i$ only through $m_0+(X-k)G_i$, so the change is exactly $-D_m+(X-k)\Delta G$. Proposition~\ref{prop:ext-customer}(b) with $dm_0=-D_m$ and a retention gain $\Delta G$ gives $dv_0=[-D_m+(v_0-k)\Delta G]/(r-g^\circ)$.

(ii) The old frontier is $m^*(g)=m_0-kg-\kappa(g-G_i)^\beta$ on $g\ge G_i$ and the new one $m_0-D_m-kg-\kappa(g-G_i-\Delta G)^\beta$ on $g\ge G_i+\Delta G$ (Proposition~\ref{prop:frontier}(a) with $G_i$ for $G$). On the new domain their difference is $d(g)=-D_m+\kappa[(g-G_i)^\beta-(g-G_i-\Delta G)^\beta]$. Since $x\mapsto x^\beta$ has an increasing derivative for $\beta>1$, $d'(g)=\beta\kappa[(g-G_i)^{\beta-1}-(g-G_i-\Delta G)^{\beta-1}]>0$, and $d(g)\to\infty$. So either $d\ge 0$ on the whole new domain (when $\kappa\Delta G^\beta\ge D_m$), or $d$ has a unique root $\hat g$, negative below and positive above it. At the root, by the mean value theorem, $D_m=\kappa\beta(\xi-G_i)^{\beta-1}\Delta G$ for some $\xi\in(\hat g-\Delta G,\hat g)$. Since $X(g)-k=\beta\kappa(g-G_i)^{\beta-1}$ (Proposition~\ref{prop:score}(d)) is increasing in $g$, the bracket follows. \qed

\paragraph{Proof of Proposition~\ref{prop:prod-breaks}.}
(a) By Proposition~\ref{prop:prod-price}(b) the cap binds if and only if $(1+\upsilon_E)\hat c_E<\hat c^\ddagger_E$, and the agent is viable if and only if $\bar p>\hat c_I$. With $\upsilon_E=0$ and the stated path, which falls strictly from $\hat c_E(0)$ toward $c_E^\infty$, the cap starts to bind when $c_E^\infty+[\hat c_E(0)-c_E^\infty]e^{-\delta_\pi T_C}=\hat c^\ddagger_E$, which gives $T_C$. (b) $\mathsf B$ has an interior inflection if and only if $(\mathfrak q-\mathfrak p)/(2\mathfrak q)\in(0,1)$, that is, $\mathfrak q>\mathfrak p$; the rest is Proposition~\ref{prop:prod-seats}(c), and AI revenue on a fixed scope at a fixed price is proportional to $\mathsf B$, so its growth in levels, $\propto\dot{\mathsf B}$, peaks at $t^*_B$. (c) restates the cited results. \qed
